\begin{table}[H]
\centering
\caption{\label{tab:manipulation-meta}Manipulation effects on the framing contrast, synthesized across surveys}
\centering
\begin{threeparttable}
\begin{tabular}[t]{lrr}
\toprule
Manipulation & Surveys & Change [95\% CI]\\
\midrule
Risky-outcome order reversed & 5 & -1.3 [-7.8, 5.3]\\
Response order reversed & 1 & 2.5 [-15.2, 20.2]\\
Risky-outcome and response order reversed & 1 & 14.1 [-2.7, 30.9]\\
Certain option fully stated (balanced) & 4 & -18.1 [-31.2, -5.0]\\
Certain option fully stated, order reversed & 1 & -17.3 [-35.3, 0.7]\\
Balanced certain option, risky-outcome order reversed & 3 & -17.5 [-65.0, 30.1]\\
Certain outcome: 199 saved & 1 & -8.3 [-26.5, 10.0]\\
Certain outcome: 201 saved & 1 & 8.5 [-9.5, 26.5]\\
Expanded drug wording & 1 & 8.7 [-6.4, 23.8]\\
\bottomrule
\end{tabular}
\begin{tablenotes}[para]
\item Note: Each entry is the change in the gain-minus-loss framing contrast, in percentage points, relative to the same survey's original wording; both arms are randomized within the survey. Rows with several surveys pool the within-survey changes with a REML random-effects model and Knapp--Hartung inference; single-survey rows report the one available difference with a normal interval. Changes from the same survey share that survey's original-wording arm, so rows are not independent of one another; within each row, surveys are independent.
\end{tablenotes}
\end{threeparttable}
\end{table}
